\section{Evaluation Methodology}
\label{sec:methodology}

% Rebuilt 2026-09-15 from docs/PAPER.md §6① and docs/experiments.md
% (comparison families, fairness, ablation matrix, measurement semantics).

Our evaluation is organized around seven questions: whether KV capacity limits
service before compute or input handling does (Q1); how many sessions each
system can serve within the objective (Q2); at what latency cost (Q3); what
each mechanism contributes (Q4); whether the resource model predicts held-out
operating points (Q5); whether state round-trips preserve retained semantics
(Q6); and where the approach stops helping (Q7).

\paragraph{Fairness.}
Formal comparisons fix the model and precision, input content and processing
semantics, actually executed generation work, reference history, scheduling
mode, batching budgets, cache and host-copy budgets, output-delivery
accounting, and observation switches. End-to-end comparisons must match work
and apply a uniform service verdict; claims that a specific mechanism causes a
benefit must control variables. Configuration names and default switches do
not substitute for verifying the executed path.

\paragraph{Compared systems.}
The comparison families are: full residency (the reference policy, and the
capacity and latency baseline), recomputation of missing state, on-demand host
reload, \systemname and its ablations, the closest session-level KV-management
systems, and retention policies such as windowing or summarization---the last
included as workload parameters that change per-session retained-state size,
not as baselines to defeat, and without adjudicating quality among retention
policies.
\placeholder{Selection and adaptation of the closest external systems
(candidates include session-KV tiering and duplex-serving systems); versions,
missing capabilities, and re-implementation status to be documented.}

\paragraph{Ablations and the three grades of timing.}
Restoration-trigger arms map to the three grades of timing information:
on-demand restoration and post-submission prefetch use no future timing;
rhythm-based prefetch uses the first two grades (rhythm and visible release
times); the release-policy arm (aligned versus staggered offsets) tests the
third grade, offset control. Reference history, host copies, and the transfer
path are held fixed across arms, and the actual moment each trigger signal is
produced is recorded. An oracle arm with exact future information, if used, is
labeled as an undeployable upper reference.

\paragraph{Metrics.}
Update service latency runs from actual submission $a(i,k)$ to generation
completion, with each session's one-time alignment wait reported separately,
and is always read together with the update-backlog metric. Capacity is
reported as offered, admitted, and completed load under the stated deadline,
distribution, and horizon. KV occupancy separates allocated space, valid
content, reusable cache, and host copies; restoration cost separates bytes,
initiation queueing, transfer, completion reporting, and rescheduling wait.
Sustained growth makes the workload non-stationary, so trajectories over time
and context are reported rather than single steady-state points.
\placeholder{Workload and platform instantiation: application timing
contracts, session-lifetime distributions, retention-policy parameters, state
geometry, and device/interconnect inventory; the matrix is not yet frozen and
no current prototype configuration is promoted to paper scope.}
